\section{Introduction}\label{sec:intro}
A real symmetric matrix can have nonnegative entries and nonnegative eigenvalues without being a sum of rank-one matrices made from nonnegative vectors. A matrix that does admit such a sum is called \emph{completely positive}. In symbols,
\[
 A=\sum_{j=1}^k u_ju_j\T,\qquad u_j\ge0.
\]
In matrix orders up to four, nonnegative entries and positive semidefiniteness are enough. From order five onward, they are not \cite{Diananda,MM}. Completely positive matrices therefore pose a natural question in matrix factorization: how can we prove that an apparently well-behaved matrix has no such representation?

There is a direct form of evidence for nonmembership. A symmetric matrix $W$ is \emph{copositive} when $x\T Wx$ is nonnegative for every nonnegative vector $x$. If $W$ is copositive but $\tr(AW)$ is negative, then $A$ cannot be completely positive: every allowed rank-one term would have nonnegative pairing with $W$. The entries of such a $W$ may be short to describe, even when finding it is difficult. A short certificate of nonmembership and a short search for that certificate are different things.

One way to look for $W$ is to follow neighboring vertices of a convex set of copositive matrices. The simplex-type method of Dutour Sikiri\'c, Sch\"urmann and Vallentin \cite{DSV} uses vertices determined by integer vectors at which a quadratic form attains its prescribed minimum. These vertices are known as \emph{perfect copositive matrices}. At each step, a pivot follows an edge to a neighboring vertex that lowers the objective $\tr(AP)$. The method may also discover a separating matrix without completing another pivot. Its appeal is that it works with local information: the current vertex, its neighboring directions, and prescribed certificate tests.

This leaves a question that a long example path does not answer. Can a choice of a different improving neighbor avoid a bad route? Could a test applied to the current direction find a certificate before the route becomes long? A lower bound for one prescribed sequence settles neither question. We study the stronger requirement that \emph{every strictly improving adjacent choice} is forced to continue and that the permitted local certificate tests fail along the entire initial part of the search.

Our first result gives a general way to construct such inputs. Begin with a rational matrix $B$ that is positive definite, has strictly positive entries, and is not completely positive. Require an explicit copositive separator and a positive description of one principal block in terms of the minimizing vectors of a suitable perfect matrix. These conditions are finite and can be certified exactly. They do not purport to characterize all nonmembers. They identify a class of starting examples, which we call \emph{certified seeds}, from which difficult local-search instances can be generated systematically.

Fix one such seed and a bounded region of initial vertices in the corresponding fixed-block face. For any integer $m\ge1$, a diagonal scaling with parameter $N$ proportional to $m+1$ produces an input requiring at least $m$ pivots for \emph{every} permitted improving choice from \emph{every} such initial vertex. The constants depend on the fixed seed and starting region. The entries of the input need only a number of binary digits proportional to $\log(m+1)$. An explicitly constructed rational copositive separator has a description of comparable length. The result remains true in a specified ball around the input. These matrices are positive definite and strictly entrywise positive, so the two elementary tests do not detect nonmembership. Thus the gap is not between the existence of a certificate and its size: it is between the existence of a short certificate and the number of local moves needed to encounter one.

The mechanism comes from a second, geometric result. Fix a positive-definite principal block $H$, and let the last column and diagonal entry of a symmetric matrix vary. Requiring its quadratic form to be at least one on nonnegative integer vectors with positive last coordinate defines a convex region. We call it a \emph{fixed-block lattice epigraph}. Its vertices are joined by edges, giving a graph in which distance counts adjacent moves. We prove, for every real positive-definite $H$, that integer changes of coordinates reduce the possible local configurations to finitely many types. In particular, the graph has bounded degree, and the number of vertices reachable within $R$ edge moves grows on the order of $R^r$ when $H$ has order $r$. This geometric statement needs neither a completely positive nonmember nor a chosen objective function. A related full-lattice version has the large-scale geometry of the integer lattice. The distances here are measured \emph{inside the fixed-block region}, which is a face when $H$ is perfect; the result makes no assertion about shortcuts outside it.

The hard inputs are not confined to the familiar five-cycle example. We give distinct graphical seed classes, as well as examples built from different perfect principal blocks. One construction produces a positive-definite, entrywise-positive nonmember of order thirteen for which \emph{every proper principal submatrix is completely positive}. Hence even checking every smaller principal submatrix exactly would not find the obstruction. We also consider stronger local tests that add diagonal matrices to combinations of previously used pivot directions. For our inputs these tests cannot separate during the forced prefix. The decomposition fact used for the resulting arrow-shaped matrices is classical \cite{SPN}; the consequence for all stored directions is part of our lower-bound result.

The simplex-type search and perfect copositive matrices were developed in earlier work \cite{DSV}; the geometry of their neighbors and related perfect-form constructions have also been studied \cite{DS,OS}. Inhomogeneous reduction proves finiteness of perfect quadratic functions up to affine integer changes of variables and scaling \cite[Section~5]{Inhom}. Here we keep a prescribed positive-definite block, include every positive last integer coordinate and the boundary imposed by nonnegativity, and prove exact intrinsic graph growth. The present amplification also adds an all-choice, finite-bit-length lower bound from a certified class of seeds. Our theorem does \emph{not} rule out algorithms that construct unrelated separators, jump between nonadjacent vertices, or choose distant starting points depending on the input. Nor does it settle eventual termination of a suitable pivot rule, as asked in \cite[Conjecture~3.2]{DSV}. The proof first establishes the fixed-block geometry, then uses the seed conditions to prevent every improving extreme direction from leaving that face too early. Concrete graphical seeds and stronger-test barriers complete the argument.
